\section{Отрицательные экраны в размерностях 5--9}\label{app:screens}

Здесь собраны подробности вычислительных кампаний, итоги которых приведены в
\S\ref{sec:open}. Все результаты этого приложения имеют статус \textbf{[Э]}:
это переборы, завершившиеся отрицательно на \emph{фиксированных} численных
формах. Они показывают, где остановились опробованные механизмы, но не
являются доказательствами невозможности --- см.\ случай индекса~$134$ ниже,
где такой барьер был обойдён сменой формы.

\subsection{Размерность 5: спуск до 132 и экраны ниже}\label{app:d5}

В $\R^5$ блочный поиск, расширение портфеля и чередование ядра с метрикой
понижают оценку АБПР с $140$ до $132$ (теорема~\ref{thm:r5}), хотя однократная
CMA-ES с бинарной целью всегда останавливалась на одном конфликте.
Индексы $139$, $136$ и $134$ были промежуточными доказанными результатами.

Специализированный конечнополевой оракул усиливает эти эвристические экраны:
после фиксации численной метрики каждая строка блока заменяется битовой маской
плохих векторов, а их произведение исчерпывается точно. Для $135$ при пороге
$0{,}978056$ проверено $1\,547\,170\,372$ набора блоков трёх типов;
для $134$ при $0{,}991322$ --- $634\,149\,671$ набор; для $133$ при
$0{,}96619$ --- $385\,308\,361$ набор: во всех случаях получена невыполнимость.
Контрольные пороги непосредственно ниже найденных значений восстанавливают
исходные ядра. На финальной численной форме промежуточной конструкции
индекса~$136$ были также исчерпаны $2\,082\,485\,047$ наборов блоков.
Смешанные блоки одной простой характеристики здесь могут перечисляться с
избытком; невыполнимость над таким надмножеством остаётся корректной. Таким образом,
дискретная оптимальность на каждой из этих \emph{фиксированных} форм проверена
глобально с точностью порядка $10^{-6}$.
Само множество плохих векторов здесь построено по метрике с плавающей точкой, поэтому это
не доказательство невозможности после изменения формы и не заменяет
рациональный сертификат теоремы. Действительно, прежний локальный потолок
$0{,}995366$ для индекса~$134$ был обойдён не шлифовкой той же формы, а
пороговым расширением портфеля на мостовой метрике: из $28$ новых точных
ядер один новый бассейн дал сертифицированное отношение
$1{,}010646405\ldots$ для промежуточного индекса~$134$. Затем прямые проверки
на один, два и три цвета ниже дали численные пересечения и для $133$, и для
$132$; последний рационализован с отношением
$1{,}010897714\ldots$. Этот пример показывает, почему невыполнимость на фиксированной метрике нельзя интерпретировать как глобальную невозможность и почему нельзя
останавливаться после проверки только $k-1$.

После сертификации 132 новый трёхшаговый экран был повторён уже относительно
этого рубежа. Для индекса~131 полный перебор всех значений с $10\,000$ рестартов,
повторный поиск на 132-метрике, деформации и доводка по активному набору сначала дали
$0{,}987424118$. Полный конечнополевой оракул затем исчерпал все
\[
 |\mathbb P^4(\mathbb F_{131})|
 =1+131+131^2+131^3+131^4=296\,765\,305
\]
проективных строк и подтвердил, что это глобальный дискретный максимум на
той фиксированной численной форме.

Следующий цикл сохранял Парето-архив различных ядер, строил общие мостовые
метрики и на каждой из них снова вызывал полный оракул. На одном мосте он
нашёл пропущенную случайными рестартами строку
\[
 \varphi(x)=x_1+52x_2+59x_3+40x_4+56x_5\pmod {131}.
\]
Четыре деформации, температурный выход и две доводки по активному набору подняли
фронтир до $0{,}990216161$. Повторное полное перечисление на новой форме
вернуло ту же строку и тот же максимум. Результат воспроизведён двумя отдельно
собранными бинарниками с разным числом потоков и разными начальными значениями генератора. Таким образом, получена
фиксированная точка точного дискретного наилучшего ответа, но не доказательство
невозможности после дальнейшего изменения формы.

Новая форма 131 рационализована независимо со знаменателями $10^4$, $10^5$ и
$10^6$; точные отношения равны соответственно $0{,}990162652$,
$0{,}990211060$ и $0{,}990215962$. Каждый независимый аудит без Qhull заново
получил $62$ фасеты, граф из $720$ вершин и $1800$ рёбер, $30$ коротких
векторов и $13$ конфликтующих пар. Во всех трёх файлах поле верхней оценки оставлено пустым.

Для $130=13\cdot5\cdot2$ две независимые многоблочные кампании и деформация
дали $0{,}976688116$, для $129=43\cdot3$ --- $0{,}957014120$.
Исчерпывающие фикс-метрические оракулы усилили эти экраны: все
$749\,112\,551$ троек строк при пороге $0{,}9766882$ оказались невыполнимыми, тогда
как $0{,}9766881$ допустим. Для 129 проверено
\[
  121\cdot3\,500\,201=423\,524\,321
\]
пара строк. Итоговые двусторонние оценки:
\[
  130:\ [0{,}9766881,0{,}9766882),\qquad
  129:\ [0{,}9570141,0{,}9570142).
\]
Эти выводы глобальны по ядрам лишь на двух
сохранённых формах и не являются доказательствами невозможности. Их
рациональные диагностики со знаменателем $10^5$ воспроизвели
$0{,}976680033$ и $0{,}957010848$, по 14 конфликтов.

\subsection{Размерность 6: кампания $343\to336$}\label{app:d6}

Для $\R^6$ оценка $\chi(\R^6)\le343$ \cite{ABPR} пока не улучшена. Точный экран
$343\to342$ перебрал $4804$ точных HNF-кандидата ($4642$ вызова полного
оракула) и достиг лишь $D/\diam=0{,}866025$; последующий многоблочный перебор всех значений и деформация
подняли фронтир для индекса~$342$ до $0{,}940393$. Проверки сразу на два и три
цвета ниже дали соответственно около $0{,}9164$ для $341$ и $0{,}935190$ для
$340$, ещё раз показывая немонотонность соседних индексов.

Более сильный эффект дала не ближайшая, а арифметически согласованная цель
$336=7\cdot4\cdot4\cdot3$. Фактор точной конструкции 343 имеет Smith-тип
$(1,1,1,7,7,7)$; мы сохранили один исходный характер над $\mathbb F_7$ и
совместно искали три остаточных блока и форму. На исходной форме лучшее ядро
имело $D/\diam=0{,}894427$. Просмотр на шаг вперёд по метрике, а затем max--min доводка по активному
набору дали ядро Smith-типа $(1,1,1,1,4,84)$ и первый фронтир
$0{,}980658185310$.

Для выхода из этого бассейна реализован точный переход через стену L-типа. Если
$\{0,v_1,\ldots,v_6\}$ --- активный унимодулярный симплекс,
$V=(v_i)$, а конкурирующая вершина $w=\sum_i\lambda_i v_i$, то зазор стены есть линейный функционал формы
\[
 \frac12\left\langle Q,\,
 ww^{\mathsf T}-\sum_i\lambda_i v_i v_i^{\mathsf T}\right\rangle.
\]
Из 97 найденных орбит стен 24 ближайшие дали 21 различную соседнюю
сигнатуру Вороного--Делоне. Строгий переход через одну из них и последующая
оптимизация по активному набору подняли отношение до $0{,}982601943196$.

Далее использован последовательный шаг SDP с отсекающими плоскостями. Для каждого
симплекса фиксированного периодического унимодулярного разбиения условие на
квадрат его радиуса имеет вид
\[
 \begin{pmatrix}
 VQV^{\mathsf T}&\operatorname{diag}(VQV^{\mathsf T})/2\\
 \operatorname{diag}(VQV^{\mathsf T})^{\mathsf T}/2&\rho
 \end{pmatrix}\succeq0.
\]
Разбиение не обязано оставаться разбиением Делоне, поэтому максимум этих радиусов всё
равно ограничивает радиус покрытия сверху. Двойственные множители проекции
короткого вектора ядра на $2\operatorname{Vor}(Q)$ одновременно дают линейную по $Q$
нижнюю оценку расстояния. Ведущая SDP-задача чередовалась с полным геометрическим
оракулом, добавлявшим пропущенные симплексы и векторы ядра. Из 720
трансляционных орбит в первой сошедшейся модели осталось 84 активных LMI и 36
двойственных сертификатов.

Два SDP-обновления и ветвление по почти активным стенам L-типа дали
\[
 D=1{,}802807715318,\qquad
 \diam=1{,}831667240842,\qquad
 D/\diam=0{,}984244122033.
\]
Численная нижняя модель на той же форме равна $0{,}984244122008$. Разрыв до
порога уменьшился с $1{,}934\%$ до $1{,}576\%$, то есть ещё на $18{,}5\%$
относительно прежнего разрыва. Это всё ещё не раскраска и не новая верхняя
оценка: SDP вычислялась в плавающей арифметике, а полный оракул подтверждает
отношение, меньшее единицы. Поэтому доказанная оценка остаётся 343.

Для массового ветвления полуопределённые конусы затем аппроксимировались
снаружи линейными спектральными сечениями
\[
  u^{\mathsf T}M(Q,\rho)u\ge0.
\]
После каждого LP-решения отрицательные собственные векторы добавлялись как
новые сечения; полученная ведущая задача решалась открытым HiGHS. На контрольной ветви
792 сечения сошлись за $0{,}022$~с, а расхождение LP-модели и полного оракула
составило около $1{,}3\cdot10^{-9}$. Для глубин $10^{-7}$, $10^{-5}$ и
$10^{-4}$ проверено по 95 ветвей по стенам L-типа (32 одиночные стены и 63 камеры
знаков). Лучшие отношения $0{,}984243977$, $0{,}984229579$ и
$0{,}984035812$ строго ухудшаются при уходе от текущей точки.

Для проверки дискретного барьера построена CP-SAT-модель строк
$[7,4,4,3]$: две строки по модулю 4 перебираются через 15 канонических
шаблонов ступенчатого вида, а мягкая цель использует целочисленные веса
$\lceil10^9(1-\rho)^4\rceil$. Счётная версия с подсказкой из предыдущей
кампании воспроизвела кандидат с двумя конфликтами и нулевым минимумом
(чистые счётные прогоны оставляли от трёх до пяти конфликтов); взвешенная
версия сохранила 27 почти допустимых конфликтов рекордного бассейна, но не
дала решения. Запуски при порогах $0{,}99$ и 1
завершились без ответа (\texttt{UNKNOWN}), а не сертификатом невозможности.
Полный HNF-архив сохранил 169 альтернативных ядер. Всем им был дан один
HiGHS-проход, затем применено последовательное отсеивание $169\to64\to24\to8$ и
дополнительная доводка 32 лидирующих или быстро растущих траекторий. Максимум
$0{,}984244122202$ отличается от прежнего значения только на
$1{,}7\cdot10^{-10}$, что меньше рабочего численного допуска, и не считается
содержательным улучшением.

Дополнительный скан сохраняющих исходный характер индексов 329, 322 и 315 после
собственной оптимизации форм дал соответственно $0{,}944303382$,
$0{,}949096819$ и $0{,}933316804$. Нециклические структуры дали
$0{,}924370435$ при $342=19\cdot3\cdot3\cdot2$,
$0{,}928954573$ при $340=17\cdot5\cdot2\cdot2$,
$0{,}908895394$ при $338=13\cdot13\cdot2$ и новый вторичный остров
$0{,}961087668$ при $333=37\cdot3\cdot3$. Это подтверждает немонотонность,
но не меняет лидера 336.

Проверена также решёточная периодическая раскраска, в которой цвет не обязан
быть одним классом смежности. Почти допустимое ядро 336 уточнялось до периода
индекса $7\cdot336=2352$, а новые слои назначались цветовым классам задачами
о совершенном паросочетании; матрица ограничений этого LP вполне
унимодулярна, поэтому релаксация имеет целочисленное оптимальное решение и
решается HiGHS. Из $19\,608$ проективных характеров 12 удаляют все петли, но для каждого
из них все шесть вариантов первого нового слоя имеют пустую строку или
столбец в графе допустимых назначений. Таким образом, этот вложенный
период не по классам смежности также не даёт 336-раскраску.

Ограничение послойного назначения затем снято полностью. Для произвольного
периода $P$ построен конечный граф Кэли
$G_P=\operatorname{Cay}(\mathbb Z^6/P,S_P)$ по точным образам всех
запрещённых сдвигов; цветом может быть любое независимое множество.
ЛП и ЦЛП покрытия множества, а также генерация столбцов
максимального веса решаются HiGHS, а массовый
оракул разрешимости следующей мощности --- CP-SAT. Поскольку граф
вершинно-транзитивен,
\[
  \chi_f(G_P)=\frac{|G_P|}{\alpha(G_P)}.
\]
Контрольная реализация на $C_5$ точно даёт $\chi_f=5/2$ и $\chi=3$.

Для простых подъёмов точного ядра 343 полностью перечислено $63$, $364$,
$3906$ и $19\,608$ остаточных профилей при $p=2,3,5,7$. При $p=2$ все графы
полные 343-дольные; при $p=3$ все 36 исключений имеют $\alpha=3$. При $p=5,7$
отдельный точный оракул на битовых масках решил все неполные графы ($3231$
при $p=5$ и $19\,536$ при $p=7$); получены соответственно $\alpha=5$ и $7$.
Тем самым исчерпывающе закрыто всё множество из $23\,941$ проективного
характера этих четырёх простых подъёмов. Составные
факторы размеров $4,8,9$ дали то же отношение $|G_P|/\alpha=343$ на
представителях всех наблюдаемых грубых профилей. Совпадение такого профиля
не доказано как изоморфизм, поэтому только составная часть этого экрана не
является исчерпывающим запретом всех исключений.

Для независимой спектральной проверки использована Fourier-редукция SDP
Ловаса абелевых графов Кэли к линейной программе по характерам
\cite{DeCorteDeLaatVallentin}. HiGHS дал $\vartheta=p$ на всех 56
представителях наблюдавшихся неполных профилей при $p=3,5,7$; максимальное
численное отклонение от целого равно $4{,}21\cdot10^{-9}$. Вместе с точным
$\alpha=p$ это по $\alpha\le\vartheta'\le\vartheta$ даёт
$\vartheta'=p$; для трёх самых разреженных представителей усиленный LP
проверен напрямую. Граница Хоффмана существенно слабее. LP-значения здесь
численные; исчерпывающий результат предыдущего абзаца обеспечен отдельным
точным оракулом.

Условие вложенности периода также снято. На классической форме $E_6^*$
точно проверены $16\,903$ различных валидных периода порядков
$686,1029,1372,1715,2401$; ни один не имеет независимого множества
необходимой мощности $3,4,5,6,8$. Отдельный скан выбранных 7-гладких
элементарно-примарных структур индексов $344$--$1026$ проверил ещё $1798$ периодов
классической формы и $1151$ период сохранённой деформированной формы кандидата
336, также без кандидата. Эти портфели случайны, не охватывают все абелевы
типы и не являются доказательством невозможности.

Проверена также принципиально иная 342-цветная семья. Пусть
$a:\mathbb Z^6\to\mathbb Z/684$ --- гомоморфизм, а один цвет объединяет два
соседних остатка. Тогда точное условие отсутствия запрещённой разности $f$ имеет
линейную целочисленную запись
\[
  2\le af-684q_f\le682.
\]
Глобальная HiGHS MIP за заданный лимит осталась без ответа, но на
лучшей деформированной форме все 15 двухкоординатных окрестностей строки
строго признаны невозможными для порога $0{,}92$. Геометрия одного цвета
теперь определяется тремя аффинными классами разностей $0,\pm1$.
Последовательность из CMA-ES, ЛП по активному набору (HiGHS) и SDP с обновлением
двойственных проекционных сертификатов стабилизировалась на
$D/\diam=0{,}9212415444$.

Чтобы буквально учесть немонотонность по периоду, для тех же 342 цветов
проверены блоки из $b=3,4,5,6$ последовательных остатков при
$N=342b$. Их точное условие равно
$b\le af-Nq_f\le N-b$. Лучший деформированный результат для $b=3$ равен
$0{,}9002398182$; полная MIP с $3921$ переменной и $3918$ ограничениями за
90 секунд прервалась по лимиту времени без ответа. Поэтому эти вычисления не
исключают всю аффинную семью, но показывают, что увеличение периода само по
себе не преодолело геометрический барьер.

Три независимых подхода --- внешняя полуопределённая релаксация, попеременная
оптимизация ядра и формы и перебор соседних арифметических типов --- упираются в
один и тот же барьер $D/\diam V_0\approx0{,}9842$. Это говорит о том, что
препятствие не в качестве локальной оптимизации: лучшей формы внутри текущего
вторичного конуса, по-видимому, нет. Поэтому индекс $336$ остаётся наиболее
перспективным кандидатом в $\R^6$, но двигаться дальше нужно, меняя сам конус и
согласуя выбор периода с выбором формы. Спектральная линия (LP по характерам в
духе Фурье и граница Хоффмана--Дельсарта) реализована и простых подъёмов не
открыла. Для ветви не по классам смежности перспективнее генерировать совместно
независимый блок остатков, строку фактора и форму, передавая двойственные веса
между ЛП и геометрическим оракулом. Спектральные методы
\cite{DSMMV} относятся к другому графу --- графу фасеточных соседей
вороновских ячеек, поэтому их оценки нельзя переносить сюда напрямую.
HiGHS остаётся естественным решателем ведущей задачи, но настоящий полуопределённый контроль
по-прежнему требует SDP-решателя и полного геометрического оракула.

\subsection{Размерности 8 и 9}\label{app:d89}

Ближайший экран $E_8$ при индексе $2400$
оставляет не менее пяти конфликтов в опробованных примарных структурах;
геометрически ориентированный кандидат начинается с $D/\diam=0{,}7071$, а
деформация за четыре поколения поднимает его лишь до $0{,}76425$. Независимый
поиск среди миллиона разреженных целочисленных соседей точной матрицы
$(3+\omega)E_8$ нашёл $722$ различных матрицы детерминанта $2400$, но лучший
фиксированный минимум равен $0{,}7071$.

Независимый экран с доводкой определителя, который точным кофактором сразу
навязывает индекс~$2400$, проверил полным оракулом ещё $540$ HNF и получил тот
же максимум $0{,}70710678$. Smith-структура
$\operatorname{diag}(1^4,7^4)$ объясняет, почему малые поправки особенно
неэффективны: для изменения определителя на единицу нужен ранг не меньше
четырёх.

Для $A_9^*$ полное запрещённое множество оказалось невыгодно
материализовывать. Сначала использовался полный ленивый оракул: для каждого
ядра перечисляются только его векторы нормы ${<}2\diam$, чего достаточно ввиду
$D(v)\ge\lVert v\rVert-\diam$. Уточнение по контрпримерам (CEGAR) с орбитами $S_{10}$ довело структуры
индексов $16\,875=3^3 5^4$ и $17\,150=2\cdot5^2 7^3$ до уровня
$0{,}81649658\ldots$.

Затем использована специальная геометрия пермутоэдра $A_9^*$. Его $1022$
релевантные фасеты строятся из тупоугольного супербазиса, а радиус покрытия
деформированной формы вычисляется прямым сканированием всех $10!$ перестановок,
без хранения $3\,628\,800$ вершин. Уточнение по контрпримерам во вторичном конусе оптимизирует форму
по архиву активных перестановок, периодически выполняет полный скан и
удерживает форму внутри конуса неравенствами тупоугольного супербазиса. Для индекса
$17\,250$ получен новый отрицательный фронтир
$D/\diam=0{,}825837653\ldots$. Независимый аудит проверил точный определитель,
все $1022$ фасеты, полный $10!$-скан и $1000$ прямых решений вершинных систем.
Для шагов $-1,-2,-3$ от $17\,253$ индексы $17\,252$, $17\,251$ и $17\,250$ дали
соответственно $0{,}656569$, $0{,}803605245$ и $0{,}825837653$;
немонотонность вновь делает проверку нескольких шагов обязательной. Это
вычислительные барьеры, не доказательства невозможности, поэтому оценки
$2401$ и $17\,253$ не изменяются.

При направленной замене фактора $71\to70$ обнаружена и исправлена смена
координат: матрица $C_9$ из \cite{ABPR} записана в базисе минимальных весов,
тогда как геометрический оракул использует фундаментальные веса. Исправленный
кандидат индекса~$17\,010$ имеет лишь $D/\diam=0{,}64196668$; прежнее значение
$0{,}627645$ было геометрически корректным, но не сохраняло заявленные пять
$\mathbb F_3$-характеров.

Следующая естественная атака --- совместное орбитно-двойственное продолжение:
расширять нарушения полными орбитами, обновлять их дуальные веса, ограничивать
дискретные изменения ядра штрафом доверительной области и одновременно делать активные
шаги по форме Грама. Для $E_8$ указанная Smith-структура даёт детерминантные
делители $343$, $49$ и $7$
для миноров порядков $7$, $6$ и $5$. Поэтому возмущение ранга меньше четырёх
меняет детерминант на число, кратное $7$, и не может дать
$2400=2401-1$; для $A_9^*$ следует
сохранять удачный $3^5$-блок Smith-разложения $3^5\cdot71$ и продолжать только
последний фактор.


